\PSsection{}{۱. اهداف کلی درس}

پس از گذراندن این درس، دانشجو قادر خواهد بود:

\begin{enumerate}
\def\labelenumi{\arabic{enumi}.}
\tightlist
\item
  مدل‌سازی عدم قطعیت در سیستم‌های مهندسی با استفاده از فضای نمونه، پیشامد و اصول موضوعه احتمال.
\item
  محاسبه احتمال شرطی، استقلال و به‌کارگیری قضیه بیز در مسایل تشخیص، ارتباطات و قابلیت اطمینان.
\item
  تحلیل متغیرهای تصادفی (یک‌متغیره و چندمتغیره) و توابع آن‌ها، امید ریاضی، گشتاورها، کوواریانس و همبستگی.
\item
  استفاده از قضیه حد مرکزی برای تقریب توزیع‌ها و تحلیل نمونه‌های بزرگ.
\item
  استنباط آماری شامل برآورد پارامتر، فاصله اطمینان و آزمون فرض.
\item
  پیاده‌سازی شبیه‌سازی مونت‌کارلو در \lr{MATLAB} برای اعتبارسنجی محاسبات نظری.
\item
  تفکر آماری برای تصمیم‌گیری مهندسی \lr{under uncertainty} (تفاوت احتمال و آمار، جامعه و نمونه، پارامتر و آماره).
\end{enumerate}

\PSsection{}{۲. جایگاه درس در برنامه درسی}

این درس پل بین \textbf{ریاضیات پایه} و \textbf{دروس تخصصی مهندسی} (ارتباطات، کنترل، پردازش سیگنال، یادگیری ماشین، قابلیت اطمینان) است.

\begin{itemize}
\tightlist
\item
  \textbf{نیمه اول (ماژول‌های ۱\lr{–}۱۱):} احتمال \lr{—} از مدل به پیش‌بینی.
\item
  \textbf{نیمه دوم (ماژول‌های ۱۲\lr{–}۱۶):} آمار \lr{—} از داده به مدل.
\end{itemize}

\PSsection{}{۳. رویکرد آموزشی}

\begin{itemize}
\tightlist
\item
  \textbf{تلفیق نظریه و شبیه‌سازی:} هر مفهوم با کد \lr{MATLAB} و آزمایش مونت‌کارلو همراه است.
\item
  \textbf{تأکید بر مثال‌های مهندسی:} کانال مخابراتی، رادار، قابلیت اطمینان، تشخیص عیب، شبکه.
\item
  \textbf{یادگیری فعال:} حل مسایل، پروژه شبیه‌سازی، تحلیل داده‌های واقعی.
\end{itemize}

\PSsection{}{۴. سرفصل تفصیلی و توزیع زمانی (۱۶ هفته)}

{\def\LTcaptype{none} % do not increment counter
\begin{longtable}[c]{@{}
  >{\centering\arraybackslash}p{(0.965\linewidth - 6\tabcolsep) * \real{0.1200}}
  >{\centering\arraybackslash}p{(0.965\linewidth - 6\tabcolsep) * \real{0.1400}}
  >{\centering\arraybackslash}p{(0.965\linewidth - 6\tabcolsep) * \real{0.3000}}
  >{\centering\arraybackslash}p{(0.965\linewidth - 6\tabcolsep) * \real{0.4400}}@{}}
\toprule\noalign{}
\rowcolor{PSnavy}\PSth{هفته} & \PSth{ماژول} & \PSth{موضوعات اصلی} & \PSth{فعالیت عملی \lr{(MATLAB)}} \\
\midrule\noalign{}
\endhead
\bottomrule\noalign{}
\endlastfoot
۱ & \lr{Module 1} & آزمایش تصادفی، فضای نمونه، پیشامد، اصول موضوعه احتمال & شبیه‌سازی سکه، تأیید قواعد جمع \\
۲ & \lr{Module 1} & احتمال شرطی، استقلال، قضیه بیز، قضیه احتمال کل & شبیه‌سازی کانال \lr{BSC}، تشخیص عیب با بیز \\
۳ & \lr{Module 2} & متغیر تصادفی یک‌متغیره، \lr{CDF}، \lr{PDF}، \lr{PMF} & ترسیم \lr{CDF/PDF} تجربی و نظری \\
۴ & \lr{Module 3} & توزیع‌های مهم (برنولی، دوجمله‌ای، پواسون، نمایی، نرمال، یکنواخت) & تولید داده تصادفی، برازش توزیع \\
۵ & \lr{Module 4} & امید ریاضی، گشتاورها، واریانس، نامساوی‌ها & تخمین گشتاورها با مونت‌کارلو \\
۶ & \lr{Module 5} & توابع متغیر تصادفی (تبدیل‌ها، توزیع مجموع، \lr{min/max}) & تبدیل متغیر و مقایسه هیستوگرام \\
۷ & \lr{Module 6} & دو متغیر تصادفی: توزیع توأم، حاشیه‌ای، شرطی & شبیه‌سازی توأم، بررسی استقلال \\
۸ & \lr{Module 7} & توابع دو متغیر تصادفی (جمع، ضرب، نسبت) & محاسبه توزیع مجموع دو \lr{RV} \\
۹ & \lr{Module 8} & امید ریاضی شرطی، پیش‌بینی، خطای \lr{MMSE} & رگرسیون خطی ساده با مونت‌کارلو \\
۱۰ & \lr{Module 9} & کوواریانس، همبستگی، ضریب همبستگی & محاسبه ماتریس کوواریانس نمونه‌ای \\
۱۱ & \lr{Module 10} & بردارهای تصادفی، توزیع نرمال چندمتغیره & تولید داده چندمتغیره نرمال \\
۱۲ & \lr{Module 11} & قضیه حد مرکزی \lr{(CLT)}، تقریب نرمال & نمایش همگرایی میانگین نمونه \\
۱۳ & \lr{Module 12} & مقدمه آمار: جامعه/نمونه، پارامتر/آماره، توزیع نمونه‌ای \lr{\ensuremath{\bar{X}}} و \lr{S\textsuperscript{2}} & شبیه‌سازی توزیع نمونه‌ای، اثر \lr{n} \\
۱۴ & \lr{Module 13} & برآورد پارامتر: \lr{MLE}، روش گشتاورها، ویژگی‌های برآوردگر & برآورد پارامتر توزیع نمایی/نرمال \\
۱۵ & \lr{Module 14} & فاصله اطمینان برای میانگین، واریانس، نسبت & ساخت \lr{CI} با شبیه‌سازی، بررسی پوشش \\
۱۶ & \lr{Module 15} & آزمون فرض: خطای نوع \lr{I/II}، توان، آزمون \lr{t}، \lr{\ensuremath{\chi}\textsuperscript{2}} & آزمون فرض با داده شبیه‌سازی‌شده \\
\end{longtable}
}

\PSsection{}{۵. روش ارزشیابی}

\begin{itemize}
\tightlist
\item
  تمرین هفتگی
\item
  آزمون میان‌ترم
\item
  آزمون پایان‌ترم
\end{itemize}
